\newpage
\section{ОСНОВНАЯ ЧАСТЬ}
\subsection{Постановка задачи}

В соответствии с выданным заданием (вариант: UID - 50430, JobName - \texttt{run\_bwa.sh}) в рамках научно-исследовательской работы необходимо решить комплекс задач, разделённый на три основных этапа:

\begin{enumerate}
    \item Подготовка виртуальной среды кластера. Необходимо развернуть и настроить виртуальную среду, имитирующую инфраструктуру распределённых вычислительных ресурсов суперкомпьютера, со следующими техническими характеристиками и требованиями:
    \begin{enumerate}
        \item Развернуть 5 виртуальных машин в среде виртуализации Oracle VM VirtualBox: один управляющий узел (\texttt{master}) с диспетчером ресурсов SLURM и средствами разработки, а также 4 вычислительных узла (\texttt{node1}--\texttt{node4}).
        \item Выделить под каждый узел конфигурацию в пределах 1--2 виртуальных ЦП и 2--4~ГБ оперативной памяти. На всех машинах должна функционировать операционная система Linux в режиме интерфейса командной строки (CLI) без графической оболочки.
        \item Настроить изолированную виртуальную сеть между всеми узлами с поддержкой разрешения имён.
        \item Обеспечить беспарольный доступ по протоколу SSH между управляющим и вычислительными узлами с авторизацией по ключам.
        \item Выполнить установку, настройку и запуск распределённого пакетного диспетчера ресурсов SLURM. Ресурсной единицей планирования определить вычислительный узел; другие типы ресурсов (ядра, оперативная память, дисковые квоты) в рамках данной работы не учитываются.
        \item Все действия в рамках экспериментов должны выполняться под единой для всех узлов непривилегированной учётной записью пользователя (\texttt{vlad}). Суперпользователь \texttt{root} применяется исключительно для конфигурирования системного ПО.
    \end{enumerate}

    \item Разработка имитатора пользовательских заявок. Программный комплекс имитации должен реализовывать статистическую модель генерации нагрузки на основе исторических данных и включать в себя:
    \begin{enumerate}
        \item Предварительный статистический анализ исторического набора данных \texttt{acct\_0921-0923}. Требуется отфильтровать записи, соответствующие индивидуальному варианту, и вычислить:
        \begin{itemize}
            \item Границы диапазона фактического времени выполнения задач ($T_{\text{min}}$ и $T_{\text{max}}$);
            \item Параметры логнормального распределения ошибки оценки времени: среднее арифметическое логарифмов ошибки ($\mu$) и их стандартное отклонение ($\sigma$).
        \end{itemize}
        \item Реализацию фазы масштабирования времени с применением понижающих коэффициентов к временным интервалам ($T_{\text{min}}$, $T_{\text{max}}$) для возможности проведения натурного эксперимента в разумные сроки без искажения математической логики модели.
        \item Использование относительной величины ошибки оценки времени $\tau = T_{\text{заявки}} / T_{\text{факт}}$ (ввиду соотношения $\Delta T > T_{\text{факт}}$ для текущего варианта) для расчёта планового времени.
        \item Алгоритм генерации отдельной заявки, где:
        \begin{itemize}
            \item Фактическое время выполнения ($t_{\text{fact}}$) моделируется как случайная величина в усечённом диапазоне $[T_{\text{min}}, T_{\text{max}}]$ по нормальному закону распределения;
            \item Запрашиваемое пользователем время ($t_{\text{req}}$) вычисляется на основе логнормального отклонения $\tau$ от сгенерированного значения $t_{\text{fact}}$;
            \item Потребность в вычислительных узлах выбирается из дискретного множества $\{1, 2, 3, 4\}$ на основе равномерного распределения.
        \end{itemize}
        \item Контроль входных параметров: имитатор должен принимать на вход объём выборки $N$ (целое положительное число, $1 \le N \le 999$).
        \item Автоматическую и последовательную отправку сформированных заявок в реальную очередь диспетчера SLURM.
    \end{enumerate}

    \item Сбор статистики и верификация эксперимента. Для оценки адекватности работы имитатора необходимо организовать систему объективного контроля:
    \begin{enumerate}
        \item Реализовать автоматизированный сбор данных непосредственно в процессе прохождения задач через диспетчер ресурсов (метрика «план -- факт»).
        \item Сохранять в независимые файлы логов как исходные сгенерированные параметры («план» имитатора), так и фактические системные метрики времени выполнения и терминальных состояний процессов, зафиксированные средствами SLURM («факт» планировщика).
    \end{enumerate}
\end{enumerate}

% \subsubsection{Подготовка виртуальной среды кластера}
% Необходимо развернуть и настроить виртуальную среду, имитирующую инфраструктуру распределённых вычислительных ресурсов суперкомпьютера, со следующими техническими характеристиками и требованиями:
% \begin{enumerate}
%     \item Развернуть 5 виртуальных машин в среде виртуализации Oracle VM VirtualBox: один управляющий узел (\texttt{master}) с диспетчером ресурсов SLURM и средствами разработки, а также 4 вычислительных узла (\texttt{node1}--\texttt{node4}).
%     \item Выделить под каждый узел конфигурацию в пределах 1--2 виртуальных ЦП и 2--4~ГБ оперативной памяти. На всех машинах должна функционировать операционная система Linux в режиме интерфейса командной строки (CLI) без графической оболочки.
%     \item Настроить изолированную виртуальную сеть между всеми узлами с поддержкой разрешения имён.
%     \item Обеспечить беспарольный доступ по протоколу SSH между управляющим и вычислительными узлами с авторизацией по ключам.
%     \item Выполнить установку, настройку и запуск распределённого пакетного диспетчера ресурсов SLURM. Ресурсной единицей планирования определить вычислительный узел; другие типы ресурсов (ядра, оперативная память, дисковые квоты) в рамках данной работы не учитываются.
%     \item Все действия в рамках экспериментов должны выполняться под единой для всех узлов непривилегированной учётной записью пользователя (\texttt{vlad}). Суперпользователь \texttt{root} применяется исключительно для конфигурирования системного ПО.
% \end{enumerate}

% \subsection*{Разработка имитатора пользовательских заявок}
% Программный комплекс имитации должен реализовывать статистическую модель генерации нагрузки на основе исторических данных и включать в себя:
% \begin{enumerate}
%     \item Предварительный статистический анализ исторического набора данных \texttt{acct\_0921-0923}. Требуется отфильтровать записи, соответствующие индивидуальному варианту, и вычислить:
%     \begin{itemize}
%         \item Границы диапазона фактического времени выполнения задач ($T_{\text{min}}$ и $T_{\text{max}}$);
%         \item Параметры логнормального распределения ошибки оценки времени: среднее арифметическое логарифмов ошибки ($\mu$) и их стандартное отклонение ($\sigma$).
%     \end{itemize}
%     \item Реализацию фазы масштабирования времени с применением понижающих коэффициентов к временным интервалам ($T_{\text{min}}$, $T_{\text{max}}$) для возможности проведения натурного эксперимента в разумные сроки без искажения математической логики модели.
%     \item Использование относительной величины ошибки оценки времени $\tau = T_{\text{заявки}} / T_{\text{факт}}$ (ввиду соотношения $\Delta T > T_{\text{факт}}$ для текущего варианта) для расчёта планового времени.
%     \item Алгоритм генерации отдельной заявки, где:
%     \begin{itemize}
%         \item Фактическое время выполнения ($t_{\text{fact}}$) моделируется как случайная величина в усечённом диапазоне $[T_{\text{min}}, T_{\text{max}}]$ по нормальному закону распределения;
%         \item Запрашиваемое пользователем время ($t_{\text{req}}$) вычисляется на основе логнормального отклонения $\tau$ от сгенерированного значения $t_{\text{fact}}$;
%         \item Потребность в вычислительных узлах выбирается из дискретного множества $\{1, 2, 3, 4\}$ на основе равномерного распределения.
%     \end{itemize}
%     \item Контроль входных параметров: имитатор должен принимать на вход объём выборки $N$ (целое положительное число, $1 \le N \le 999$).
%     \item Автоматическую и последовательную отправку сформированных заявок в реальную очередь диспетчера SLURM.
% \end{enumerate}

% \subsection*{Сбор статистики и верификация эксперимента}
% Для оценки адекватности работы имитатора необходимо организовать систему объективного контроля:
% \begin{enumerate}
%     \item Реализовать автоматизированный сбор данных непосредственно в процессе прохождения задач через диспетчер ресурсов (метрика «план -- факт»).
%     \item Сохранять в независимые файлы логов как исходные сгенерированные параметры («план» имитатора), так и фактические системные метрики времени выполнения и терминальных состояний процессов, зафиксированные средствами SLURM («факт» планировщика).
% \end{enumerate}
% \vfill